\documentclass{article}
\usepackage{graphicx, physics, dsfont}
\usepackage{graphicx}
\usepackage{fancyhdr}
\usepackage{hyperref}
\usepackage{ragged2e}
\usepackage{amsmath, amsthm}
\usepackage[margin=1in]{geometry}
\usepackage{setspace}
\usepackage{tikz}
\usetikzlibrary{positioning}
\usepackage[
]{geometry}
\pagestyle{fancy}
\lhead{Zoeb Izzi}
\rhead{\today}

\rfoot{}
\begin{document}
\thispagestyle{fancy}
\begin{center}\LARGE{Global Quantum Mechanics Challenge \\ Solution Set}\end{center}
\section{Problem A}
Region (a) is \textbf{radio} waves, which encompass about \textbf{100 km to 1 mm} in wavelength.\\
Region (b) is \textbf{infrared} waves, which encompass about\textbf{ 1mm to 700 nm }in wavelength.\\
Region (c) is \textbf{ultraviolet} waves, which encompass about \textbf{400 nm to 10 nm} in wavelength.\\
Region (d) is \textbf{X-rays}, which encompass about \textbf{10 nm to 0.01 nm} in wavelength.\\
Region (e) is \textbf{gamma} waves, which encompass wavelengths \textbf{less than 0.01 nm}.
\section{Problem B}
We know that:
\[E = \frac{h c}{\lambda}\]
Thus, for a \textbf{red photon} with wavelength 650 nm, we have:
\[E = \frac{hc}{650 \cdot 10^{-9}} = 3.05607 \cdot 10^{-19} \text{ J}= \boxed{1.90745 \text{ eV}}\]
And, for a \textbf{green photon} with wavelength 532 nm, we have:
\[E = \frac{hc}{650 \cdot 10^{-9}} = 3.73302 \cdot 10^{-19} \text{ J} = \boxed{2.33053 \text{ eV}}\]
\section{Problem C}
\subsection{Part A}
We know that the photon just needs to have that work function amount of energy in order to excite the electron into leaving. If we use the above formula, 
\[E = \frac{h c}{\lambda}, \lambda = \frac{hc}{E}\]
so therefore,
\[\lambda = \frac{hc}{2.3 \text{ eV}} = \frac{hc}{3.68501\cdot10^{-19} \text{ J}} = 5.3906 \cdot 10^{-7} \text{ m} \approx \boxed{539 \text{ nm}}\]
is the minimum wavelength needed.
\subsection{Part B}
\[2500 \text{ Å} = 250 \text{ nm} = \lambda\]
\[E = \frac{hc}{\lambda} = \frac{1240}{250} \text{ eV} = \boxed{4.96 \text{ eV } = 7.9468\cdot 10^{-19} \text{ J}}\]
\section{Problem D}
\subsection{Part A}
\[\Delta x \Delta p \ge\frac{h}{4\pi}\]
\[\Delta x \frac{h}{\lambda} \ge \frac{h}{4\pi}\]
\[\Delta x \ge \frac{\lambda}{4\pi}\]
\[\boxed{\Delta x \ge }\]
\subsection{Part B}
Here, $p = mv = 75 \cdot 3 \text{ Nm} = 225 \text{ Nm}$
\[\Delta x \ge \frac{h}{4\pi \cdot 225}\]
\[\boxed{\Delta x \ge }\]
\section{Problem E}
\subsection{Quantum Computing Uses}
Scientists hope quantum computers will be especially useful in state-
\subsection{Classical Computing Advantages}
Classical computers are still better than quantum computers for most tasks because operating systems need to be highly stable, which makes the probabilistic nature of quantum mechanics problematic. 
Also, the overhead provided by quantum computers is (for the time being), in most cases, too high to justify their use in anything but highly specialized scenarios (such as factoring large numbers or finding elements in unsorted lists).

\end{document}